\documentclass[11pt,a4paper,twocolumn]{article}
\usepackage[a4paper,margin=2cm,columnsep=0.7cm]{geometry}
\usepackage{fontspec}
\setmainfont{texgyretermes}[Path=fonts/, Extension=.otf, UprightFont=*-regular, BoldFont=*-bold, ItalicFont=*-italic, BoldItalicFont=*-bolditalic]
\usepackage{xeCJK}
\setCJKmainfont{NotoSerifCJKsc}[Path=fonts/, Extension=.otf, UprightFont=*-Regular, BoldFont=*-Bold]
\xeCJKsetup{AutoFakeSlant=true}
\usepackage{amsmath}
\usepackage{booktabs}
\usepackage[table]{xcolor}
\usepackage{tikz}
\usepackage{titlesec}
\usepackage{ragged2e}
\definecolor{accent}{HTML}{000000}
\titleformat{\section}{\color{accent}\bfseries\large}{\thesection.}{0.5em}{}
\titlespacing*{\section}{0pt}{1.4ex plus .2ex}{0.3ex}
\setlength{\parindent}{0pt}
\setlength{\parskip}{0.5em}
\title{\vspace{-1.5em}\bfseries\color{accent}\Large 个人理财系统的基石：三大支柱}
\author{}
\date{}

\begin{document}

\twocolumn[
  \begin{@twocolumnfalse}
    \maketitle
    \vspace{-3.8em}
    \begin{center}
      \small Nonthawit Doungsodsri\\
      nonthawit.com\\
      2026年6月3日\\
      版本 1.0.0
    \end{center}
    \vspace{0.5em}
    \begin{quote}
      \normalsize\textbf{摘要} —— 本文提出一套简洁而持久的个人理财框架，即“三大支柱体系”：现金流管理（Cashflow）、投资与储蓄（价值保值）。其核心主张在于：长期财务成功并非源于复杂的策略，而在于\textbf{纪律}，以及持续遵循自身体系五年乃至十年以上的坚守。
    \end{quote}
    \vspace{1em}
  \end{@twocolumnfalse}
]

\section{引言}
个人理财是在自身人生目标的指引下，对个人资金进行管理的实践。对大多数人而言，问题并非收入不足，而是缺乏一套清晰的\textit{体系}来管理已有的收入。本文提供一套贯穿一生、切实可用的体系，其根基在于一个核心等式。

\begin{center}
{\setlength{\fboxsep}{10pt}\fbox{$\displaystyle \text{收入} \;-\; \text{支出} \;=\; \text{自由基金}$}}
\end{center}

“自由基金”是通往财务自由的原材料。若结果为负，首要目标便是将其转为正值；若已为正，下一个问题便是如何分配。本文的答案是：将自由基金引导至协同运作的三大支柱之中——现金流管理、投资与储蓄。每根支柱各司其职，以下将逐一阐述。

\section{支柱一 —— 现金流管理（Cashflow）}
\textbf{定义。}现金流管理，是指通过持续记录收支，掌握每月资金的流入与流出状况。无需复杂工具，一本简单的笔记本或一款应用程序即已足够。关键在于持之以恒，而非追求精密。

\textbf{目标。}每月清楚知晓真实的“自由基金”——这个数字告诉你拥有多少可用于投资与储蓄的余力。没有它，一切计划都不过是猜测。这根支柱是其他两根支柱赖以站立的\textit{根基}。

\textbf{实践。}坚持每日或每周记账，直至形成习惯；月末复盘收支总览；区分必要开支与自由支配开支；并以将“自由基金”稳固保持在正值为先，再迈向第二与第三支柱。

\section{支柱二 —— 投资}
\textbf{定义。}投资，是指让金钱为你工作——将自由基金导入能够产生长期回报的资产之中。与持有闲置现金不同，经过投资的资金有机会借助复利（compounding）随时间增值。

\textbf{目标。}获取超越单纯储蓄的回报，令财富加速增长，使你更接近财务自由。核心原则是：\textbf{只投资自己真正了解的领域}。跟风投资、不理解所持资产，只会为不可控的风险敞开大门。

\textbf{实践。}切勿操之过急；投资如同植树，无法催促其生长。先深入学习，待理解透彻后，再运用\textbf{DCA}（Dollar-Cost Averaging）：无论市场涨跌，每月定额买入资产。此法将情绪从决策中剥离，塑造持久的投资纪律，尤其适合初学者。当然，DCA只是诸多入市方式之一。其他方式同样存在，例如一次性投入（lump-sum）、再平衡（rebalancing）或价值投资（value investing），各有优势，适合不同情境。读者应深入研究，选择契合自身目标与性情的方式。

\section{支柱三 —— 储蓄（价值保值）}
\textbf{定义。}此处的储蓄，并非指积累闲置现金，而是持有\textbf{能够保值的资产}，以降低生活中的风险。闲置的现金每年都在因通货膨胀（inflation）而侵蚀购买力；良好的储蓄，必须守护资金价值不受损耗。

\textbf{目标。}建立安全感与抵御风险的缓冲，使意外来临时，手中有能够保值的资产可以依靠。这根支柱承担“防守”职责，而投资支柱则承担“进攻”职责。

\textbf{实践。}与其执着于某一单一资产，不如依据\textbf{保值的特质}加以甄选，例如：抗通胀能力、合理的流动性以及广泛的认可度。当下可参考的示例包括黄金、指数基金（index fund）、债券（bond）或用于转移风险的保险，而这些选择会随时代变迁而更迭。因此，读者应研究当时可用的选项，分散配置以避免风险过度集中于单一资产，并像对待投资一样，维持稳定的储蓄比例。

\section{分配与纪律}
理解三大支柱之后，实践层面的问题便是：如何将“自由基金”在三者之间进行分配。三大支柱并非彼此割裂，而是协同支撑同一目标——财务自由，如图~\ref{fig:pillars} 所示。

\begin{figure}[t]
\centering
\begin{tikzpicture}[font=\small]
  \fill[accent!18] (-0.2,2.1) rectangle (6.4,2.5);
  \draw[accent] (-0.2,2.1) rectangle (6.4,2.5);
  \node[accent] at (3.1,2.3) {\bfseries 财务自由};
  \foreach \x/\lbl in {0/Cashflow, 2.2/Investment, 4.4/Savings} {
    \fill[accent!15] (\x,0) rectangle (\x+1.6,2.0);
    \draw[accent] (\x,0) rectangle (\x+1.6,2.0);
    \node[accent,align=center] at (\x+0.8,1.0) {\bfseries \lbl};
  }
  \draw[accent,line width=1pt] (-0.4,0) -- (6.6,0);
\end{tikzpicture}
\caption{支撑财务自由的三大支柱。}
\label{fig:pillars}
\end{figure}

\textbf{分配示例。}假设月收入 5{,}000 元，必要支出 3{,}000 元，自由基金为 2{,}000 元。可将这笔自由基金按 30/40/30 的比例进行分配，如表~\ref{tab:alloc} 所示。需要特别强调的是：30/40/30 仅为说明性示例，并非固定公式。每个人的收入、义务、目标与风险承受能力各不相同，适合他人的比例未必适合你。读者应根据自身实际情况设计专属比例。唯一不可省略的是：自由基金必须始终流入\textbf{投资}与\textbf{储蓄}两根支柱——二者是必选项，任何方案都必须包含。至于比例大小，以及其余部分（如自由支配或其他目标）如何分配，读者可自由设计，以契合各自的生活。

\begin{table}[t]
\centering
\small
\begin{tabular}{lrr}
\toprule
\rowcolor{accent!15}
\textbf{去向} & \textbf{比例} & \textbf{金额（元）} \\
\midrule
投资   & 30\% & 600 \\
储蓄   & 40\% & 800 \\
自由支配 & 30\% & 600 \\
\midrule
\textbf{自由基金合计} & \textbf{100\%} & \textbf{2{,}000} \\
\bottomrule
\end{tabular}
\caption{2,000 元自由基金的分配示例（仅供说明）。}
\label{tab:alloc}
\end{table}

\textbf{纪律是核心。}无论如何分配，决定成败的因素既非比例，也非复杂策略，而是\textbf{持续遵循自身体系五年乃至十年以上}。DCA 与持续储蓄的力量，只有在你让时间与复利（compounding）充分发挥作用时，方能显现。拥有一套体系的最深远价值在于：它能消弭情绪的干扰——无论是牛市中的贪婪，还是波动中的恐惧——这才是长期理财的真正敌人。遵循体系，决策便依据原则而非一时的感受。凡是早起步、持之以恒者，终将胜过那些聪明却三天打鱼两天晒网的人。

\section{结论}
构建稳健的个人财务，无需高深的金融知识，也无需秘诀策略，只需一套易于理解的体系，以及切实可行的纪律。三大支柱——以现金流管理掌握自由基金、以投资令金钱为己所用、以储蓄降低风险——协同运作，持之以恒，将一路提升你财务生活的质量。

请铭记这一原则：\textbf{“从简单的体系开始，然后坚持去做。”}时间与坚持，是每一位投资者最强大的盟友。

\vspace{1em}
\noindent\rule{\linewidth}{0.4pt}\\
{\small\raggedright 本文档由 \textbf{Nonthawit Doungsodsri} 创作 —— 为公共利益发布。更多内容见 nonthawit.com\par}

\end{document}
